\chapter{反馈、锁存器与触发器}

组合逻辑没有记忆。按钮按下时输出为 1，按钮松开后输出回到 0。若硬件需求是“按一下后保持亮”，只用组合逻辑就不够了。要保存过去，电路必须让当前输出影响未来输出，也就是把输出的一部分反馈回输入。

\section{一、反馈为什么能保存一个 bit}

两个反相器首尾相接，会形成两个稳定状态。若左边节点为 0，右边经过反相后为 1；右边的 1 再反相回来，正好维持左边为 0。反过来，左边为 1、右边为 0 也能稳定。只要没有外部强制改变，这个状态会一直保持。

\begin{pdfdiagram}[x=1cm,y=1cm]
  \node[hw compute,minimum width=1.4cm] (inv1) at (3.0,1.5) {反相器};
  \node[hw compute,minimum width=1.4cm] (inv2) at (6.3,1.5) {反相器};
  \node[hw state,minimum width=0.95cm] (q) at (1.2,1.5) {Q};
  \node[hw state,minimum width=0.95cm] (nq) at (8.1,1.5) {$\overline{Q}$};
  \draw[hw bus] (q.east) -- (inv1.west);
  \draw[hw bus] (inv1.east) -- (inv2.west);
  \draw[hw bus] (inv2.east) -- (nq.west);
  \draw[BookTeal,line width=0.55pt,-{Latex[length=2mm,width=1.5mm]}] (nq.south) .. controls (6.4,0.15) and (2.7,0.15) .. (q.south);
  \node[hw label,align=center] at (4.65,0.45) {反馈让输出反过来维持输入\\形成两个可稳定停留的状态};
  \node[hw label,text=BookAccent] at (3.0,2.25) {状态 1：Q=0};
  \node[hw label,text=BookAccent] at (6.35,2.25) {$\overline{Q}$=1};
\end{pdfdiagram}
\diagramnote{两个反相器相互支撑，电路不是记住“变量名”，而是停在一个稳定物理状态上。}

这件事说明“保存一个 bit”不是抽象魔法，而是电路中某些节点在反馈作用下停留在一个稳定状态。一个稳定状态解释为 0，另一个稳定状态解释为 1。状态的本质，是电路有能力在输入不再变化时仍然维持某个内部结果。

但单纯反馈还不够。若我们无法可靠地把它从 0 改成 1，或从 1 改成 0，它就只是一个会停住的环，不是可用存储单元。因此下一步要加入写入控制。

\section{二、锁存器：允许时跟随，禁止时保持}

锁存器可以理解为带写使能的 1-bit 存储单元。写使能有效时，输入 D 可以影响输出 Q；写使能无效时，反馈通路维持原来的 Q。

\begin{longtable}{p{0.22\linewidth}p{0.24\linewidth}p{0.42\linewidth}}
\toprule
写使能 & 数据输入 & 输出行为 \\
\midrule
有效 & 0 或 1 & 输出跟随输入变化 \\
无效 & 任意 & 输出保持原值 \\
\bottomrule
\end{longtable}

例如一个“按住才能设置”的电路中，enable 为 1 时，D=1 会把锁存器写成 1；enable 回到 0 后，即使 D 变回 0，Q 仍保持 1。这样，过去某一刻的输入就被硬件保存下来了。

锁存器的问题在于透明窗口。只要 enable 有效，输入变化就可能一路穿透到输出。若多个锁存器共用同一个 enable，并且它们之间还接着组合逻辑，状态可能在一个打开窗口内穿过多级，边界变得模糊。

\section{三、触发器把写入压到时钟边界}

触发器解决的是“什么时候写入”的问题。它不像锁存器那样在一段时间内透明，而是在时钟沿附近采样输入，把采样值保存到输出。

\begin{pdfdiagram}[x=1cm,y=1cm]
  \node[hw label,text=BookNavy] at (2.2,2.55) {锁存器：enable 有效时透明};
  \draw[BookRule,line width=0.55pt] (0.2,1.85) -- (4.2,1.85);
  \draw[BookAccent,line width=0.75pt] (0.2,1.35) -- (1.1,1.35) -- (1.1,1.85) -- (3.0,1.85) -- (3.0,1.35) -- (4.2,1.35);
  \node[hw label] at (-0.2,1.35) {en};
  \draw[BookTeal,line width=0.75pt] (0.2,0.65) -- (0.9,0.65) -- (0.9,1.05) -- (1.65,1.05) -- (1.65,0.55) -- (2.45,0.55) -- (2.45,0.95) -- (3.7,0.95);
  \node[hw label] at (-0.2,0.75) {D};
  \draw[BookRed,line width=0.75pt] (0.2,0.1) -- (1.1,0.1) -- (1.1,1.05) -- (1.65,1.05) -- (1.65,0.55) -- (2.45,0.55) -- (2.45,0.95) -- (3.0,0.95) -- (3.0,0.1) -- (4.2,0.1);
  \node[hw label] at (-0.2,0.1) {Q};
  \node[hw label,align=center] at (2.2,-0.55) {enable 打开期间\\D 的变化可能穿到 Q};

  \node[hw label,text=BookNavy] at (8.2,2.55) {触发器：只在边界采样};
  \draw[BookRule,line width=0.55pt] (6.2,1.85) -- (10.2,1.85);
  \draw[BookAccent,line width=0.75pt] (6.2,1.35) -- (6.95,1.35) -- (6.95,1.85) -- (7.0,1.35) -- (8.05,1.35) -- (8.05,1.85) -- (8.1,1.35) -- (9.15,1.35) -- (9.15,1.85) -- (9.2,1.35) -- (10.2,1.35);
  \node[hw label] at (5.8,1.35) {clk};
  \draw[BookTeal,line width=0.75pt] (6.2,0.65) -- (6.9,0.65) -- (6.9,1.05) -- (7.7,1.05) -- (7.7,0.55) -- (8.65,0.55) -- (8.65,0.95) -- (10.2,0.95);
  \node[hw label] at (5.8,0.75) {D};
  \draw[BookRed,line width=0.75pt] (6.2,0.1) -- (6.95,0.1) -- (6.95,0.65) -- (8.05,0.65) -- (8.05,1.05) -- (9.15,1.05) -- (9.15,0.55) -- (10.2,0.55);
  \node[hw label] at (5.8,0.1) {Q};
  \node[hw label,align=center] at (8.2,-0.55) {只有时钟沿附近\\Q 才接收 D};
\end{pdfdiagram}
\diagramnote{锁存器在窗口内跟随输入，触发器把写入压到离散时钟边界。}

\begin{lstlisting}
on clock edge:
  Q = D
between clock edges:
  Q holds
\end{lstlisting}

这让大电路有了清楚节奏。一个时钟沿之后，上一批触发器输出新状态；这些状态经过组合逻辑传播，形成下一批输入；到下一个时钟沿，目标触发器再统一采样。于是“计算”和“保存”被分开了。

可以用一个两步计数需求理解它。第一步，当前值进入加一电路；第二步，下一个时钟沿把加一结果保存。如果没有触发器边界，加一结果又会立刻回到输入，再被加一，再回到输入，电路会在一个周期内不断变化。触发器把反馈切成一拍一拍的更新。

\section{四、亚稳态来自真实采样边界}

触发器不是数学上的瞬间拍照。若输入 D 正好在时钟沿附近变化，内部节点可能短时间落在不稳定区域，既不像明确 0，也不像明确 1。这个现象叫亚稳态。

亚稳态不是因为设计者不会写逻辑表达式，而是物理器件在边界条件下真实存在的行为。多数情况下，它会在一小段时间后落回 0 或 1，但落回哪边、需要多久，无法按普通组合逻辑那样精确保证。

这直接带来设计纪律：被触发器采样的输入，必须在时钟沿前后保持稳定。来自另一个时钟节奏或外部世界的信号，不能直接拿来控制大面积电路，通常要先经过同步结构，降低亚稳态扩散的概率。

\section{五、状态边界把计算切成步骤}

本章得到三件关键东西。反馈提供两个可稳定停留的状态；锁存器提供“允许写入、否则保持”的能力；触发器把写入限制到时钟边界。

从这里开始，硬件不再只是“输入一变，输出跟着变”的组合网络，而可以分步骤工作。下一章会把这些状态边界组织成全局节奏：时钟周期、复位、建立时间、保持时间和关键路径。

\begin{classicnote}
本章第一次出现真正的“机器状态”。状态不是变量名，而是由反馈和存储节点保持住的一组 bit。

\begin{itemize}
\item 纸上实验：画出一个 D 触发器前后各接一段组合逻辑，标出当前周期的旧 Q、组合路径上的 D、下一个时钟沿后的新 Q。
\item 状态证据：锁存器在使能窗口内跟随输入，触发器只在时钟边界采样；这两个边界不能混用。
\item 亚稳态证据：若 D 在采样窗口附近变化，不能用普通真值表预测内部节点何时稳定，只能用同步器降低风险。
\item 检查题：一个反馈加一电路如果没有触发器，为什么可能在一个周期内连续变化？触发器到底切断了哪条路径？
\end{itemize}

经典教材通常会把“组合逻辑”和“状态逻辑”分开建模：组合逻辑计算下一值，状态元件在边界保存下一值。后面 CPU 的 PC、寄存器、状态机和控制器都遵守这个模型。
\end{classicnote}
